% =============================================================================
% CAPÍTULO 12: IoT, Telemetria e o Gateway SUS-Twin
% Volume 2 — Gêmeos Digitais Periodontais
% V2 Standalone — sem \input{}
% Seções:
%   12.1 - Arquitetura IoT para Periodontia Digital
%   12.2 - Gateway SUS-Twin: Processamento de Borda com MQTT
%   12.3 - Dashboard Clínico com Grafana e InfluxDB
%   12.4 - Projeto Prático: Monitoramento Oclusal IoT Completo
%   12.5 - Protocolos IoT e Interoperabilidade Odontológica
%   12.6 - Exercícios Propostos
% =============================================================================
\label{ch:plataformas}

% --- BADGE DO CAPÍTULO ---
\begin{highlight}
\large\textbf{Badge do Capítulo}\\[4pt]
\rule{\textwidth}{1pt}\\[4pt]
\begin{tabular}{ll}
\textbf{Nível-alvo:} & N2 --- Pesquisa Reprodutível \\
\textbf{Habilidades:} & configurar MQTT, gateway de borda, \\
                     & dashboard IoT em tempo real \\
\textbf{Pré-requisitos:} & Python 3.11+, paho-mqtt, Docker \\
\textbf{Comando de verificação:} & \texttt{tdd\_validate.py --chapter 12}
\end{tabular}
\end{highlight}

\destaque{Nível-alvo N2 --- Pesquisa Reprodutível: ao final deste capítulo, o leitor será capaz de implementar o gateway MQTT do Framework SUS-Twin para coleta de telemetria de sensores intraorais, processar dados na borda e visualizá-los em dashboard clínico.}

% =============================================================================
\section{Arquitetura IoT para Periodontia Digital}
% =============================================================================

A Internet das Coisas\footnote{Internet das Coisas (IoT): rede de dispositivos físicos equipados com sensores, software e conectividade de rede que coletam e trocam dados. Na odontologia, sensores intraorais IoT permitem monitoramento contínuo de variáveis biomecânicas e fisiológicas.} (IoT) aplicada à periodontia digital constitui a camada sensorial do Framework SUS-Twin. Sem sensores confiáveis e uma arquitetura de comunicação robusta, o gêmeo digital periodontal é apenas uma simulação desconectada da realidade clínica. A IoT odontológica viabiliza a coleta sistemática de grandezas físicas e químicas diretamente do ambiente bucal do paciente, alimentando o modelo preditivo com dados reais e contínuos.

Os sensores intraorais\footnote{Sensor: dispositivo que converte uma grandeza física (força, temperatura, pH) em um sinal elétrico mensurável. Sensores intraorais são miniaturizados para uso em próteses, alinhadores ou dispositivos retentivos.} contemporâneos operam em três domínios principais: biomecânico (força oclusal, torque, contato), térmico (temperatura intraoral) e químico (pH salivar, biomarcadores inflamatórios). A Tabela~\ref{tab:sensores-grandezas} mapeia os tipos de sensor, as grandezas mensuradas e exemplos de aplicação clínica em periodontia.

\begin{table}[hbt!]
\centering
\caption{Sensores IoT intraorais, grandezas mensuradas e aplicações periodontais.}
\label{tab:sensores-grandezas}
\small
\begin{tabular}{lllp{5.5cm}}
\toprule
\textbf{Tipo de Sensor} & \textbf{Grandeza} & \textbf{Exemplo Comercial} & \textbf{Aplicação Periodontal} \\
\midrule
Strain gauge piezoelétrico & Força oclusal (N) & T-Scan Novus (50--1000\,N) & Mapeamento de contatos prematuros em próteses sobre implantes \\
\addlinespace
Termopar tipo K & Temperatura (°C) & Thermochron iButton (25--50\,°C) & Monitoramento de inflamação gengival pós-cirúrgica \\
\addlinespace
ISFET & pH salivar & Sensirion SPS30 (0--14) & Risco de cárie e desmineralização peri-implante \\
\addlinespace
Acelerômetro MEMS & Movimento mandibular (mm/s) & MPU-6050 (3 eixos) & Análise de bruxismo e apertamento noturno \\
\addlinespace
Eletrodo seco & EMG superficial ($\mu$V) & MyoWare 2.0 & Atividade dos masseteres e temporais na ATM \\
\bottomrule
\end{tabular}
\end{table}

A telemetria\footnote{Telemetria: processo de coleta e transmissão remota de dados de sensores para uma central de processamento. Na odontologia, permite que o clínico acompanhe a evolução do paciente sem consultas presenciais frequentes.} odontológica impõe requisitos específicos de arquitetura: (a) baixa latência para detecção de eventos oclusais (tipicamente $<100$\,ms), (b) capacidade de operação off-line em clínicas com conectividade intermitente, e (c) segurança de dados compatível com a LGPD. O gateway SUS-Twin atende a esses três requisitos por meio de uma arquitetura de borda --- edge computing\footnote{Edge computing (computação de borda): processamento de dados realizado próximo à fonte geradora (o sensor), em vez de em servidores remotos. Reduz latência, economiza banda e preserva privacidade.} ---, onde o processamento preliminar ocorre no próprio gateway local antes da transmissão para a nuvem clínica.

% =============================================================================
\section{Gateway SUS-Twin: Processamento de Borda com MQTT}
% =============================================================================

O gateway\footnote{Gateway: dispositivo ou software que faz a ponte entre sensores locais e a infraestrutura de rede, traduzindo protocolos, agregando dados e executando processamento de borda.} SUS-Twin é o núcleo da arquitetura IoT do framework. Implementado em Python 3.11+, ele utiliza o protocolo MQTT\footnote{MQTT (Message Queuing Telemetry Transport): protocolo de mensagens publish/subscribe leve e eficiente, projetado para IoT. Opera sobre TCP/IP e requer um broker central para distribuir as mensagens entre dispositivos.} (Message Queuing Telemetry Transport) como espinha dorsal de comunicação. MQTT foi escolhido por sua eficiência em banda estreita, suporte nativo a qualidade de serviço (QoS) e modelo publish/subscribe que desacopla sensores de consumidores de dados.

A Figura conceitual do gateway segue o fluxo: sensores $\rightarrow$ broker MQTT $\rightarrow$ processador de borda $\rightarrow$ InfluxDB $\rightarrow$ Grafana. O broker\footnote{Broker MQTT: servidor central que recebe mensagens publicadas por clientes (sensores) e as distribui para clientes inscritos (assinantes). Mosquitto é o broker open-source mais utilizado.} Mosquitto gerencia o roteamento das mensagens entre os componentes, permitindo que múltiplos sensores publiquem em tópicos hierárquicos como \texttt{sustwin/patient/\{id\}/occlusal}.

\paragraph{Código 1: Publicador de sensor intraoral simulado.}

\begin{lstlisting}[style=pythonstyle, caption=Publicador MQTT que simula sensor intraoral T-Sense Pro, label=cod:publisher]
#!/usr/bin/env python3
"""publicador_sensor.py — Simula sensor intraoral SUS-Twin."""
import paho.mqtt.client as mqtt
import json, time, random
from datetime import datetime

# Configuracoes do broker
BROKER = "localhost"
PORT = 1883
TOPIC = "sustwin/patient/001/occlusal"
CLIENT_ID = "tsense_pro_001"

def on_connect(client, userdata, flags, rc):
    print(f"[PUB] Conectado ao broker (codigo {rc})")

def on_publish(client, userdata, mid):
    print(f"[PUB] Mensagem {mid} publicada com sucesso")

client = mqtt.Client(client_id=CLIENT_ID)
client.on_connect = on_connect
client.on_publish = on_publish
client.connect(BROKER, PORT, 60)
client.loop_start()

try:
    while True:
        # Simula leitura do sensor intraoral
        payload = {
            "timestamp": datetime.now().isoformat(),
            "patient_id": "001",
            "sensor_id": CLIENT_ID,
            "force_occlusal_N": round(random.uniform(50, 500), 2),
            "temperature_C": round(random.uniform(33, 38), 1),
            "ph_saliva": round(random.uniform(6.2, 7.6), 2),
            "occlusal_contact_ms": round(random.uniform(10, 250), 1),
            "battery_pct": round(random.uniform(60, 100), 1)
        }
        result = client.publish(TOPIC, json.dumps(payload), qos=1)
        time.sleep(5)  # Publica a cada 5 segundos
except KeyboardInterrupt:
    print("[PUB] Encerrando publicador...")
    client.loop_stop()
    client.disconnect()
\end{lstlisting}

\paragraph{Código 2: Processador de borda com detecção de anomalias.}

O processador de borda --- edge processor --- assina o tópico MQTT e executa análises em janela deslizante antes de persistir os dados no banco temporal.

\begin{lstlisting}[style=pythonstyle, caption=Assinante MQTT com processamento de borda e detecção de anomalias, label=cod:edge]
#!/usr/bin/env python3
"""edge_processor.py — Processamento de borda do gateway SUS-Twin."""
import paho.mqtt.client as mqtt
import json
from collections import deque
import statistics

WINDOW_SIZE = 10  # Janela deslizante de 10 amostras
data_window = deque(maxlen=WINDOW_SIZE)

def on_message(client, userdata, msg):
    """Callback executado a cada mensagem recebida."""
    payload = json.loads(msg.payload)
    data_window.append(payload)

    if len(data_window) >= WINDOW_SIZE:
        forces    = [d["force_occlusal_N"] for d in data_window]
        temps     = [d["temperature_C"]    for d in data_window]
        contacts  = [d["occlusal_contact_ms"] for d in data_window]

        # Estatisticas de borda
        stats = {
            "window_size":      WINDOW_SIZE,
            "mean_force_N":     round(statistics.mean(forces), 2),
            "max_force_N":      round(max(forces), 2),
            "std_force_N":      round(statistics.stdev(forces), 2),
            "mean_temp_C":      round(statistics.mean(temps), 2),
            "contact_imbalance": max(contacts) - min(contacts),
            "anomaly":          max(forces) > 450 or max(temps) > 37.5
        }

        print(json.dumps(stats, indent=2))

        if stats["anomaly"]:
            print("[ALERTA] Anomalia oclusal detectada!")
            print("[ALERTA] Encaminhar para revisao clinica imediata.")

# Configuracao do assinante
client = mqtt.Client(client_id="edge_gateway_001")
client.on_message = on_message
client.connect("localhost", 1883, 60)
# Assina todos os pacientes (wildcard +)
client.subscribe("sustwin/patient/+/occlusal", qos=1)
print("[EDGE] Gateway SUS-Twin aguardando dados...")
client.loop_forever()
\end{lstlisting}

A detecção de anomalias na borda --- como forças oclusais acima de 450\,N ou temperatura intraoral superior a 37,5\,°C --- dispara alertas imediatos sem depender de conectividade com a nuvem, garantindo capacidade de resposta mesmo em clínicas com internet instável. Este é o princípio fundamental do design do gateway SUS-Twin: processar o máximo possível na borda para minimizar latência e dependência de infraestrutura externa.

% =============================================================================
\section{Dashboard Clínico com Grafana e InfluxDB}
% =============================================================================

O pipeline de dados completo do gateway SUS-Twin persiste as séries temporais no InfluxDB\footnote{InfluxDB: banco de dados especializado em séries temporais (time-series database), otimizado para armazenar e consultar dados de sensores ao longo do tempo. Utiliza a linguagem de consulta Flux.} e as visualiza no Grafana\footnote{Grafana: plataforma open-source de visualização e monitoramento que cria dashboards interativos a partir de múltiplas fontes de dados, incluindo InfluxDB, Prometheus e PostgreSQL.}, formando o dashboard clínico em tempo real.

A configuração do InfluxDB pode ser realizada via Docker com o comando:

\begin{lstlisting}[style=bashstyle, caption=Inicialização do InfluxDB e Mosquitto via Docker Compose, label=cod:docker]
# docker-compose.yml — Infraestrutura IoT do SUS-Twin
version: "3.8"
services:
  mosquitto:
    image: eclipse-mosquitto:2
    ports:
      - "1883:1883"
      - "9001:9001"
    volumes:
      - ./mosquitto.conf:/mosquitto/config/mosquitto.conf

  influxdb:
    image: influxdb:2.7
    ports:
      - "8086:8086"
    environment:
      DOCKER_INFLUXDB_INIT_MODE: setup
      DOCKER_INFLUXDB_INIT_USERNAME: admin
      DOCKER_INFLUXDB_INIT_PASSWORD: sustwin2026
      DOCKER_INFLUXDB_INIT_ORG: sustwin
      DOCKER_INFLUXDB_INIT_BUCKET: telemetria_oclusal

  grafana:
    image: grafana/grafana:11.0
    ports:
      - "3000:3000"
    depends_on:
      - influxdb
\end{lstlisting}

Após a inicialização dos containers, o código a seguir escreve os dados processados no InfluxDB:

\begin{lstlisting}[style=pythonstyle, caption=Escrita de dados processados no InfluxDB, label=cod:influx]
#!/usr/bin/env python3
"""influxdb_writer.py — Persiste telemetria no InfluxDB."""
from influxdb_client import InfluxDBClient, Point
from influxdb_client.client.write_api import SYNCHRONOUS
import json, time

URL     = "http://localhost:8086"
TOKEN   = "sustwin2026"
ORG     = "sustwin"
BUCKET  = "telemetria_oclusal"

client  = InfluxDBClient(url=URL, token=TOKEN, org=ORG)
write_api = client.write_api(write_type=SYNCHRONOUS)

def persist_telemetry(data: dict) -> None:
    """Converte dict em Point do InfluxDB e persiste."""
    point = Point("oclusal_reading") \
        .tag("patient_id",  data.get("patient_id", "unknown")) \
        .tag("sensor_id",   data.get("sensor_id", "unknown")) \
        .field("force_N",           float(data["force_occlusal_N"])) \
        .field("temperature_C",     float(data["temperature_C"])) \
        .field("ph_saliva",         float(data["ph_saliva"])) \
        .field("contact_ms",        float(data["occlusal_contact_ms"])) \
        .field("battery_pct",       float(data["battery_pct"]))
    write_api.write(bucket=BUCKET, record=point)

# Exemplo de uso (dados simulados)
if __name__ == "__main__":
    sample = {
        "timestamp": "2026-06-21T14:30:00",
        "patient_id": "001",
        "sensor_id": "tsense_pro_001",
        "force_occlusal_N": 235.7,
        "temperature_C": 35.2,
        "ph_saliva": 6.8,
        "occlusal_contact_ms": 120.5,
        "battery_pct": 85.0
    }
    persist_telemetry(sample)
    print("Dado persistido no InfluxDB com sucesso.")
\end{lstlisting}

O dashboard\footnote{Dashboard: painel visual que consolida indicadores-chave em tempo real, permitindo ao clínico monitorar múltiplos pacientes simultaneamente a partir de gráficos, tabelas e alertas coloridos.} Grafana pode ser configurado com consultas Flux para visualizar a evolução da força oclusal ao longo do tempo, médias diárias e detecção de padrões anômalos. Um exemplo de consulta Flux para o painel de força oclusal é:

\begin{lstlisting}[style=pythonstyle, caption=Consulta Flux para dashboard Grafana, label=cod:flux]
from(bucket: "telemetria_oclusal")
  |> range(start: -24h)
  |> filter(fn: (r) => r._measurement == "oclusal_reading")
  |> filter(fn: (r) => r._field == "force_N")
  |> aggregateWindow(every: 5m, fn: mean)
  |> yield(name: "mean_force_5min")
\end{lstlisting}

% =============================================================================
\section{Projeto Prático: Monitoramento Oclusal IoT Completo}
% =============================================================================

Este projeto prático integra todos os componentes do capítulo em um cenário clínico realístico. O objetivo é implementar o monitoramento oclusal contínuo de um paciente com prótese total sobre implantes, utilizando o gateway SUS-Twin.

\paragraph{Cenário clínico.} Paciente 40 anos, sexo feminino, portadora de prótese total fixa inferior (protocolo Branemark). Apresenta histórico de peri-implante e queixa de desconforto oclusal no lado direito. O periodontista solicita monitoramento IoT por 48 horas para mapear os padrões de força oclusal antes de ajuste oclusal definitivo.

\paragraph{Infraestrutura necessária.}
\begin{itemize}
    \item Docker Desktop (Windows) ou Docker Engine (Linux)
    \item Python 3.11+ com pacotes: \texttt{paho-mqtt}, \texttt{influxdb-client}, \texttt{requests}
    \item Navegador web para acessar Grafana (porta 3000) e InfluxDB UI (porta 8086)
\end{itemize}

\paragraph{Passos para reprodução.}

\begin{enumerate}
    \item \textbf{Inicie a infraestrutura:} execute \texttt{docker compose up -d} no diretório com o \texttt{docker-compose.yml} do Código~\ref{cod:docker}. Verifique se Mosquitto (porta 1883), InfluxDB (8086) e Grafana (3000) estão ativos.
    \item \textbf{Execute o publicador:} em um terminal, rode \texttt{python publicador\_sensor.py} (Código~\ref{cod:publisher}). O sensor simulado publicará leituras oclusais no tópico MQTT a cada 5 segundos.
    \item \textbf{Execute o processador de borda:} em outro terminal, rode \texttt{python edge\_processor.py} (Código~\ref{cod:edge}). Observe as estatísticas da janela deslizante e os alertas de anomalia.
    \item \textbf{Persista os dados:} execute \texttt{python influxdb\_writer.py} (Código~\ref{cod:influx}) para registrar leituras no banco temporal.
    \item \textbf{Configure o Grafana:} acesse \url{http://localhost:3000} (admin/admin), adicione InfluxDB como fonte de dados, importe a consulta Flux do Código~\ref{cod:flux} e crie um painel com:
    \begin{itemize}
        \item Gráfico de séries temporais: força oclusal (N) nas últimas 24h
        \item Medidor de temperatura intraoral (°C)
        \item Tabela de alertas de anomalia
        \item Estatísticas descritivas (média, máximo, desvio padrão)
    \end{itemize}
\end{enumerate}

\paragraph{Validação.} O projeto é considerado concluído quando o dashboard Grafana exibir dados de telemetria atualizados em tempo real, com detecção visual de picos de força oclusal e alertas configurados.

% =============================================================================
\section{Protocolos IoT e Interoperabilidade Odontológica}
% =============================================================================

A escolha do protocolo de comunicação IoT impacta diretamente a latência, o consumo energético e a confiabilidade do sistema de telemetria. A Tabela~\ref{tab:protocolos-iot} compara os principais protocolos disponíveis e sua adequação ao contexto odontológico.

\begin{table}[hbt!]
\centering
\caption{Protocolos IoT e sua aplicação em odontologia digital.}
\label{tab:protocolos-iot}
\small
\begin{tabular}{lp{2.8cm}p{1.8cm}p{4cm}p{2.5cm}}
\toprule
\textbf{Protocolo} & \textbf{Transporte} & \textbf{Padrão} & \textbf{Uso Odontológico} & \textbf{Limitações} \\
\midrule
\textbf{MQTT} & TCP/IP & Publish/Subscribe & Telemetria oclusal, sensores intraorais contínuos & Requer broker central \\
\addlinespace
\textbf{CoAP} & UDP & Request/Response & Sensores de baixo consumo (pilha botão) & Confiabilidade reduzida em rede instável \\
\addlinespace
\textbf{HTTP/2} & TCP/IP & Request/Response & Sincronização em nuvem, prontuário eletrônico & Alto overhead, não recomendado para tempo real \\
\addlinespace
\textbf{LwM2M} & UDP/CoAP & Gerenciamento & Provisionamento remoto de sensores clínicos & Complexidade de implementação \\
\addlinespace
\textbf{WebSocket} & TCP/IP & Full-duplex & Dashboard em tempo real no navegador & Consumo de memória elevado \\
\bottomrule
\end{tabular}
\end{table}

Para o Framework SUS-Twin, o MQTT com QoS nível 1 (garantia de entrega pelo menos uma vez) é o protocolo recomendado para a camada de sensoriamento contínuo, enquanto HTTP/2 é utilizado para a sincronização periódica com o prontuário eletrônico do paciente (e-SUS Odonto / HL7 FHIR).

% =============================================================================
\section{Exercícios Propostos}
% =============================================================================

Os exercícios a seguir consolidam os conceitos práticos do capítulo, progredindo da configuração básica até a integração completa do pipeline IoT.

\paragraph{Exercício 1: Publicação de dado MQTT simulado.}
Implemente um script Python que publique, a cada 10 segundos, dados simulados de um sensor de temperatura intraoral no tópico \texttt{sustwin/patient/002/temperature}. O payload deve conter os campos: \texttt{timestamp}, \texttt{patient\_id}, \texttt{sensor\_id}, \texttt{temperature\_C}, \texttt{battery\_pct}. Utilize o broker público \texttt{test.mosquitto.org} (porta 1883) para testar. Inclua tratamento de exceção para falhas de conexão.

\paragraph{Exercício 2: Configuração de broker Mosquitto local.}
Crie um arquivo \texttt{mosquitto.conf} completo para um broker local com as seguintes configurações:
\begin{itemize}
    \item Porta 1883 (MQTT padrão) e 9001 (WebSocket)
    \item Autenticação por usuário/senha (arquivo \texttt{password.txt})
    \item ACL (Access Control List) restringindo o tópico \texttt{sustwin/} apenas para usuários autenticados
    \item Persistência de mensagens em arquivo (\texttt{persistence true})
    \item Log nível \texttt{information}
\end{itemize}
Inicie o broker com \texttt{docker run} e valide a conexão com \texttt{mosquitto\_sub} e \texttt{mosquitto\_pub}.

\paragraph{Exercício 3: Criação de dashboard Grafana completo.}
Utilizando os dados persistidos no InfluxDB pelo Código~\ref{cod:influx}, crie um dashboard Grafana com os seguintes painéis:
\begin{enumerate}
    \item Série temporal da força oclusal (N) com média móvel de 15 minutos
    \item Histograma de distribuição das forças oclusais nas últimas 6 horas
    \item Tabela de alerta listando todos os eventos com força $>400$\,N
    \item Medidor em gauge (\textit{single stat}) da temperatura intraoral atual
\end{enumerate}
Exporte o dashboard como JSON e salve em \texttt{dashboards/sustwin-oclusal.json}. O comando de verificação \texttt{tdd\_validate.py --chapter 12} deve ser capaz de importar e validar o JSON exportado.
